\documentclass[a4paper,10pt]{article}

% Настройка шрифтов и языка (для XeLaTeX)
\usepackage{fontspec}
\usepackage{xunicode}
\usepackage{xltxtra}
\setmainfont{Liberation Serif} % Или другой шрифт с поддержкой кириллицы
\usepackage[russian]{babel}

% Пакеты для оформления
\usepackage[margin=1.5cm, top=2cm, bottom=2cm]{geometry}
\usepackage{xcolor}
\usepackage{tabularx}
\usepackage{colortbl}
\usepackage{enumitem}
\usepackage{fancyhdr}
\usepackage{amssymb}
\usepackage{titlesec}
\usepackage{needspace}


% Определение цветов
\definecolor{sputnikBlue}{RGB}{31, 73, 125}
\definecolor{sputnikRed}{RGB}{192, 0, 0}
\definecolor{tableHeader}{RGB}{54, 95, 145}

%для гиперссылок
\usepackage[colorlinks=true, linkcolor=sputnikBlue, urlcolor=sputnikBlue]{hyperref}

% Настройка заголовков
\titleformat{\section}{\color{sputnikBlue}\normalfont\Large\bfseries}{}{0pt}{\thesection. }
\titleformat{\subsection}{\color{sputnikBlue}\normalfont\large\bfseries}{}{0pt}{\thesubsection\ }

% Команды для мастей
\newcommand{\clb}{\textcolor{black}{♣}}
\newcommand{\dia}{\textcolor{red}{♦}}
\newcommand{\hrt}{\textcolor{red}{♥}}
\newcommand{\spd}{\textcolor{black}{♠}}

\newcommand{\NT}{БК}

\begin{document}

% Шапка документа
\begin{itemize}[label=\textcolor{sputnikRed}{\tiny\(\blacksquare\)}, leftmargin=*, noitemsep]
    \item \textbf{Баланс на гейм} – 24+ очка. Геймовые контракты – 3\NT, 4\hrt/\spd, 5\clb/\dia. Минорный гейм играем на 9-картном фите 
    \item \textbf{Фит} – наличие на одной линии 8+ карт в какой-либо масти.
    \item \textbf{Открытие мастью на 1 уровне} обещает от 12 очков. С 6+ очками в ответ НЕ ПАСУЕМ.
    \item \textbf{Вход мастью} – 8-15 очков 5+ карт на первом уровне, 10-15 очков 5+ карт на втором уровне без скачка, или по правилам блоков (можно со скачком). входы на 3+ уровне как чувствуем или тоже по правилам блоков
    \item \textbf{Вход 1\NT(2\NT после блока оппонентов)} – 16-19 очков (можно и с 15), равномерная рука, есть держка в масти оппонента.
\end{itemize}

\section{Спектр открытий}
\begin{itemize}[label=\textcolor{sputnikRed}{\textbullet}, leftmargin=*, noitemsep]
    \item \textcolor{sputnikRed}{Выбранное открытие отрицает все остальные открытия.} % Матвей: ну с этим я бы поспорил, но да ладно
    \item \textcolor{sputnikRed}{Выбранное открытие в масти отрицает 5-картную масть старше.}
\end{itemize}

\noindent
\begin{tabularx}{\textwidth}{|l|l|X|}
\hline
\rowcolor{tableHeader} \textcolor{white}{\textbf{Открытие}} & \textcolor{white}{\textbf{Очков}} & \textcolor{white}{\textbf{Значение}} \\ \hline
1\clb & 12-21 & 2+\clb\ , нет 5 карт в \hrt/\spd\, нет 4 карт в \dia/ : \\
      &       & Равномерный расклад только при 12-14 и 18-19 очках. \\ \hline
1\dia & 12-21 & 4+\dia, нет 5 карт в \hrt/\spd\ : \\
      &       & Равномерный расклад только при 12-14 или 18-19 очках. \\ \hline
1\hrt & 12-21 & 5+\hrt, нет 5 карт в \spd \\ \hline
1\spd & 12-21 & 5+\spd \\ \hline
1\NT  & 15-17 & Расклад 5332 / 4333 / 4432, нет 5 карт в \hrt/\spd \\ \hline
2\clb & 22+ (либо 9 взяток с руки)  & Любой расклад, ФГ \\ \hline
2\dia/\hrt/\spd & 5-10 & 6 карт в масти открытия, не менее половины очков в масти \\ \hline
2\NT  & 20-21 & Расклад 5332 / 4333 / 4432, нет 5 карт в \hrt/\spd \\ \hline
3\clb & 5-9 & 7\clb, почти все очки в \clb ИЛИ очень хорошие 6\clb \\ \hline
3\dia/\hrt/\spd & 5-9 & 7 карт в масти открытия, почти все очки в масти \\ \hline
4\clb/\dia/\hrt/\spd & 5-9 & 8 карт в масти открытия, почти все очки в масти \\ \hline
\end{tabularx}

\section{Ответы на открытие 1\clb}
\noindent
\begin{tabularx}{\textwidth}{|X|X|}
\hline
\rowcolor{tableHeader} \textcolor{white}{\textbf{Отвечающий}} & \textcolor{white}{\textbf{Ответ открывающегося}} \\ \hline
1\dia\ = 6+ очков, 4+\hrt, нет 5\spd, Ф1. & \begin{itemize}[label=\textcolor{sputnikRed}{\tiny\(\blacksquare\)}, leftmargin=*, nosep]
    \item 1\hrt = 12+, ровно 3 карты в \hrt
    \item 2\hrt = 12-14, ровно 4 карты в \hrt
    \item 3\hrt = 15-17, ровно 4 карты в \hrt
    \item 4\hrt = 18-21, ровно 4 карты в \hrt
    \item 1\spd = 12+, ровно 4 карты в \spd, нет 3 карт в \hrt
    \item 2\clb = 6+ карт, 12-14, нет 3 карт в \hrt и 4 карт в \spd
    \item 3\clb = 6+ карт, 15-17, нет 3 карт в \hrt и 4 карт в \spd
    \item 1\NT = любой другой расклад, в том числе 18+ очков (\hyperlink{sec:double-checkback}{!задаем двойной чекбек!})
\end{itemize} \\ \hline
1\hrt\ = 6+ очков, 4+\spd, Ф1. & \begin{itemize}[label=\textcolor{sputnikRed}{\tiny\(\blacksquare\)}, leftmargin=*, nosep]
    \item 1\spd = 12+, ровно 3 карты в \spd
    \item 2\spd = 12-14, ровно 4 карты в \spd
    \item 3\spd = 15-17, ровно 4 карты в \spd
    \item 4\spd = 18+, ровно 4 карты в \spd
    \item 2\clb = 6+ карт, 12-14, нет 3 карт в \spd
    \item 3\clb = 6+ карт, 15-17, нет 3 карт в \spd
    \item 1\NT = любой другой расклад, в том числе 18+ очков (\hyperlink{sec:double-checkback}{!задаем двойной чекбек!})
\end{itemize} \\ \hline
1\NT\ = 6-9 очков, нет 4 карт в \hrt/\spd, НФ &
\begin{itemize}[label=\textcolor{sputnikRed}{\tiny\(\blacksquare\)}, leftmargin=*, nosep]
    \item 2\clb = 6+ карт, 12-14
    \item 2\clb = 6+ карт, 15-17
    \item 2\NT = инвит в 3бк
    \item 3\NT = контракт который хотим играть
\end{itemize} \\ \hline
2\NT\ = 10-11 очков, нет 4 карт в \hrt/\spd, инвит в 3\NT, НФ &
\begin{itemize}[label=\textcolor{sputnikRed}{\tiny\(\blacksquare\)}, leftmargin=*, nosep]
    \item 3\NT = инвит принят
    \item пас = инвит не принят
\end{itemize} \\ \hline
3\NT\ = 12 очков, нет 4 карт в \hrt/\spd, НФ &
\begin{itemize}[label=\textcolor{sputnikRed}{\tiny\(\blacksquare\)}, leftmargin=*, nosep]
    \item пас
\end{itemize} \\ \hline

\end{tabularx}

\section{Ответы на открытие 1\dia}
\noindent
\begin{tabularx}{\textwidth}{|X|X|}
\hline
\rowcolor{tableHeader} \textcolor{white}{\textbf{Отвечающий}} & \textcolor{white}{\textbf{Ответ открывающегося}} \\ \hline
1\hrt\ = 6+ очков, 4+\hrt, нет 5\spd, Ф1. & \begin{itemize}[label=\textcolor{sputnikRed}{\tiny\(\blacksquare\)}, leftmargin=*, nosep]
    \item 2\hrt = 12-14, ровно 4 карты в \hrt
    \item 3\hrt = 15-17, ровно 4 карты в \hrt
    \item 4\hrt = 18-21, ровно 4 карты в \hrt
    \item 1\spd = 12+, ровно 4 карты в \spd, нет 4 карт в \hrt
    \item 2\clb = 5-4 \clb\dia/\dia\clb, 12-14, нет 4 карт в \hrt/\spd
    \item 3\clb = 5-4 \clb\dia/\dia\clb, 15-17, нет 4 карт в \hrt/\spd
    \item 2\dia = 6+ \dia, 12-14, нет 4 карт в \hrt/\spd
    \item 3\dia = 6+ \dia, 15-17, нет 4 карт в \hrt/\spd
    \item 1\NT = любой другой расклад, в том числе 18+ очков (\hyperlink{sec:double-checkback}{!задаем двойной чекбек!})
\end{itemize} \\ \hline

1\spd\ = 6+ очков, 4+\spd, Ф1. & \begin{itemize}[label=\textcolor{sputnikRed}{\tiny\(\blacksquare\)}, leftmargin=*, nosep]
    \item 2\spd = 12-14, ровно 4 карты в \spd
    \item 3\spd = 15-17, ровно 4 карты в \spd
    \item 4\spd = 18+, ровно 4 карты в \spd
    \item 2\dia = 6+ \dia, 12-14, нет 4 карт в \spd
    \item 3\dia = 6+ \dia, 15-17, нет 4 карт в \spd
    \item 2\clb = 5-4 \clb\dia/\dia\clb, 12-14, нет 4 карт в \hrt/\spd
    \item 3\clb = 5-4 \clb\dia/\dia\clb, 15-17, нет 4 карт в \hrt/\spd
    \item 1\NT = любой другой расклад, в том числе 18+ очков (\hyperlink{sec:double-checkback}{!задаем двойной чекбек!})
\end{itemize} \\ \hline
1\NT\ = 6-9 очков, нет 4 карт в \hrt/\spd, НФ &
\begin{itemize}[label=\textcolor{sputnikRed}{\tiny\(\blacksquare\)}, leftmargin=*, nosep]
    \item 2\dia = 6+ \dia, 12-14
    \item 2\dia = 6+ \dia, 15-17
    \item 2\clb = 5-4 \clb\dia/\dia\clb, 12-14, нет 4 карт в \hrt/\spd
    \item 3\clb = 5-4 \clb\dia/\dia\clb, 15-17, нет 4 карт в \hrt/\spd
    \item 2\NT = инвит в 3бк
    \item 3\NT = контракт который хотим играть
\end{itemize} \\ \hline
2\NT\ = 10-11 очков, нет 4 карт в \hrt/\spd, инвит в 3\NT, НФ &
\begin{itemize}[label=\textcolor{sputnikRed}{\tiny\(\blacksquare\)}, leftmargin=*, nosep]
    \item 3\NT = инвит принят
    \item пас = инвит не принят
\end{itemize} \\ \hline
3\NT\ = 12 очков, нет 4 карт в \hrt/\spd, НФ &
\begin{itemize}[label=\textcolor{sputnikRed}{\tiny\(\blacksquare\)}, leftmargin=*, nosep]
    \item пас
\end{itemize} \\ \hline

\end{tabularx}

\section{Ответы на открытие 1\hrt\ (на 1\spd\ аналогично)}
\noindent
\begin{tabularx}{\textwidth}{|X|X|X|}
\hline
\rowcolor{tableHeader} \textcolor{white}{\textbf{Отвечающий}} & \textcolor{white}{\textbf{Открывающий}} & \textcolor{white}{\textbf{Отвечающий}} \\ \hline
1\spd\ на 1\hrt\ = 6+ очков, 4+\spd, Ф1 & Фит 2/3/4 ур., 1\NT, повтор \hrt, 2\clb/\dia\ (5+\hrt\ и 4+\clb/\dia) & \\ \hline
2\hrt\ на 1\hrt\ =6-9 очков, ровно 3 карты в\hrt & 
\begin{itemize}[label=\textcolor{sputnikRed}{\tiny\(\blacksquare\)}, leftmargin=*, nosep]
    \item ПАС=12-14
    \item 3\hrt\ =15-17
    \item 4\hrt\ =18+
\end{itemize} & \\ \hline
1\NT\ = 6-11 очков без фита ИЛИ 10-11 очков с ровно 3 картами в \hrt\ , нет 4\spd, Ф1 & 
\begin{itemize}[label=\textcolor{sputnikRed}{\tiny\(\blacksquare\)}, leftmargin=*, nosep]
    \item 2\NT = 15-17 очков, 5 карт в \hrt
    \item 3\NT = 18+ очков, 5 карт в \hrt
    \item 2\dia = 12-14 очков, 5 карт в \hrt, 4 карты в \dia
    \item 2\clb = \textbf{12-14 очков, 5 карт в \hrt}
    \item 2\hrt = 12-14 очков, 6+ карт в \hrt
    \item 3\hrt = 15+ очков, 6+ карт в \hrt
\end{itemize} &
\begin{itemize}[label=\textcolor{sputnikRed}{\tiny\(\blacksquare\)}, leftmargin=*, nosep]
    \item Масть на ур.2 =слабая рука, 6 карт
    \item Масть на ур.3 =инвитная рука, 6 карт
    \item 2\NT = натур.инвит
    \item 3\NT =гейм.рука 4-3-3-3 (с фитом, контракт на выбор)
    \item 2\hrt = дуплет, к игре
    \item 3\hrt = 10-11 очков с 3-картным фитом
\end{itemize} \\ \hline 
2\clb\ =12+ очков, равномер/фит в мажоре/ 5+ \clb\ , нет 5 карт в \dia/\hrt , ФГ &
\begin{itemize}[label=\textcolor{sputnikRed}{\tiny\(\blacksquare\)}, leftmargin=*, nosep]
    \item Повтор мажора открытия=6 карт
    \item Новая масть= 4+ карт
    \item 2\NT = расклад 5332
\end{itemize}
& \\ \hline
2\dia\ =12+ очков,5+ \dia, ФГ & 
\begin{itemize}[label=\textcolor{sputnikRed}{\tiny\(\blacksquare\)}, leftmargin=*, nosep]
    \item Повтор мажора открытия=6 карт
    \item Новая масть= 4+ карт
    \item 2\NT = расклад 5332
\end{itemize}
& \\ \hline
2\hrt на 1\spd =12+ очков,5+ \hrt , ФГ & 
\begin{itemize}[label=\textcolor{sputnikRed}{\tiny\(\blacksquare\)}, leftmargin=*, nosep]
    \item Фит 3+ карты
    \item Повтор масти открытия=6+ карт
    \item Новая масть= 4+ карт
    \item 2\NT = расклад 5332
\end{itemize} & \\ \hline
2\spd\ на 1\hrt =12+ очков, 6+ \spd , ФГ & 
\begin{itemize}[label=\textcolor{sputnikRed}{\tiny\(\blacksquare\)}, leftmargin=*, nosep]
    \item Фит 2+ карты
    \item Повтор масти открытия=6+ карт
    \item Новая масть= 4+ карт
\end{itemize}& \\ \hline
2\NT\ = 10-11 очков с 4-картным фитом или фитом с двумя фигурами &  & \\ \hline
3\hrt\ =блокируещего характера с 4-картным фитом &  & \\ \hline
\end{tabularx}

\section{Ответы на открытие 1\NT}
\noindent
\begin{tabularx}{\textwidth}{|X|X|X|}
\hline
\rowcolor{tableHeader} \textcolor{white}{\textbf{Отвечающий}} & \textcolor{white}{\textbf{Открывающий}} & \textcolor{white}{\textbf{Отвечающий}} \\ \hline
2\clb\ (\textbf{Стейман}) = 0+ очков (но надо четко знать что делать на каждый вариант событий), есть мажорная четверка & 
\begin{itemize}[label=\textcolor{sputnikRed}{\tiny\(\blacksquare\)}, leftmargin=*, nosep]
    \item 2\dia = нет мажорных 4-ок
    \item 2\hrt = 4 карты в \hrt, могут быть 4 карты в \spd
    \item 2\spd = 4 карты в \spd , нет 4 карт в \hrt
\end{itemize}&
НА ОТВЕТ 2\dia, если 4-5 или 5-5 в мажорах:
\begin{itemize}[label=\textcolor{sputnikRed}{\tiny\(\blacksquare\)}, leftmargin=*, nosep]
    \item 2\hrt= слабая рука с 5-5, открывающий выбирает козырь
    \item 2\spd = инвитная рука с 5-картной \spd----при 5-5, если теперь открывающий поставил 2\NT, ставим 3\hrt
    \item 3\spd = геймовая рука с 5-картной \spd----при 5-5, если теперь открывающий поставил 3\NT, ставим 4\hrt
\end{itemize}
\\ \hline
2\dia\ (\textbf{трансфер в \hrt}) = 0+ очков, 5+\hrt & 2\hrt\ (ОБЯЗАТЕЛЬНЫЙ ПРИЕМ) &
\begin{itemize}[label=\textcolor{sputnikRed}{\tiny\(\blacksquare\)}, leftmargin=*, nosep]
    \item Пас=0-6 очков
    \item 2\spd = инвит, 5 черв и 4 пик
    \item 3\spd = гейм, 5 черв и 4 пик
    \item 3 \hrt = 6+ карт, 7-8 очков
    \item 4 \hrt = 6+ карт, 9+ очков
    \item Новая масть=4+ карт, 10+ очков, Фг (Открывающий после ставит 4\hrt, если фит был, или торговля дальше)
    \item 2\NT = расклад 5332, 8-9 очков (C 17 очками Открывающий после ставит 4\hrt, если фит был, или 3\NT, если нет)
    \item 3\NT = расклад 5332, 10+ очков, Гейм (Открывающий после переставит в 4\hrt, если фит был)
\end{itemize}\\ \hline
2\hrt\ (\textbf{трансфер в \spd}) = 0+ очков, 5+\spd & 2\spd\ (ОБЯЗАТЕЛЬНЫЙ ПРИЕМ) & Аналогично трансферу в \hrt \\ \hline
2\spd\ (\textbf{трансфер в \clb}) = 0+ очков (См.далее), 6+\clb & 3\clb\ (ОБЯЗАТЕЛЬНЫЙ ПРИЕМ) & 
\begin{itemize}[label=\textcolor{sputnikRed}{\tiny\(\blacksquare\)}, leftmargin=*, nosep]
    \item Пас=слабая рука
    \item ФГ=еще заявка
\end{itemize}
\\ \hline
3\clb\ (\textbf{трансфер в \dia}) = 0+ очков (См.далее), 6+\dia & 3\dia\ (ОБЯЗАТЕЛЬНЫЙ ПРИЕМ) & 
\begin{itemize}[label=\textcolor{sputnikRed}{\tiny\(\blacksquare\)}, leftmargin=*, nosep]
    \item Пас=слабая рука
    \item ФГ=еще заявка
\end{itemize}
\\ \hline
2\NT\ = 8-9 очков, инвит к 3\NT & 3\NT\ с максимумом, иначе ПАС & \\ \hline
4\clb\ = 15+ очков, вопрос о количестве тузов & 
\begin{itemize}[label=\textcolor{sputnikRed}{\tiny\(\blacksquare\)}, leftmargin=*, nosep]
    \item 4\dia = 1 туз
    \item 4\hrt = 2 туза
    \item 4\spd = 3 туза
    \item 4\NT = 0 или 4 туза
\end{itemize} & \\ \hline
\end{tabularx}

\section{Ответы на открытие 2\clb\ Форсинг гейм}
\noindent
\begin{tabularx}{\textwidth}{|X|X|}
\hline
\rowcolor{tableHeader} \textcolor{white}{\textbf{Отвечающий}} & \textcolor{white}{\textbf{Открывающий}} \\ \hline
2\dia\ = 0--8, любая рука &
\begin{itemize}[label=\textcolor{sputnikRed}{\tiny\(\blacksquare\)}, leftmargin=*, nosep]
    \item 2\NT\ равномер, далее торговля как после открытия 2БК (стейман, трансферы)
    \item 2\hrt/\spd/\dia/\clb\ при 5+
\end{itemize} \\ \hline
2\hrt/\spd\ = 9+ очков, 5+ \hrt/\spd & Согласование масти или 2\NT/3\NT \\ \hline
2\NT\ = 9+ очков, нет 5 в \hrt/\spd & Согласование масти или 2\NT/3\NT \\ \hline
3\clb/\dia\ = 9+ очков, 5+ \clb/\dia & Согласование масти или 2\NT/3\NT \\ \hline
\end{tabularx}

\section{Двойной чекбек}
\hypertarget{sec:double-checkback}{}
\subsection{Определение чекбэковой позиции, важные моменты}
\textbf{Чекбэковая позиция-позиция, в которой были сделаны нами 3 заявки на уровне 1.}
\subsection{Конвенция}
\noindent
\begin{tabularx}{\textwidth}{|X|X|}
\hline
\rowcolor{tableHeader} \textcolor{white}{\textbf{Отвечающий}} & \textcolor{white}{\textbf{Открывающий}} \\ \hline
2\clb=просьба сказать 2\dia:

1)Намерение играть 2\dia

2)=сила инвита&
\begin{itemize}[label=\textcolor{sputnikRed}{\tiny\(\blacksquare\)}, leftmargin=*, nosep]
    \item 2\dia=принятие. !ДАЛЕЕ ЕСЛИ У ОТВЕЧАЮЩЕГО СЛУЧАЙ 2, то на 2\dia делаем натуральную заявку!
    \item натуральная иная заявка=много очков, нельзя останавливаться в 2\dia.
\end{itemize} \\ \hline
2\dia = произвольный ФГ & Натурально \\ \hline
2\hrt/\spd/ = слабая рука, натурально & ... \\ \hline
2\NT = просьба сказать 3\clb:

1)слабая рука с \clb

2)шлемовая рука с \clb & 3\clb, 

!ОТВЕЧАЮЩИЙ ПРОДОЛЖАЕТ ТОРГОВЛЮ со ШЛЕМОВОЙ рукой!\\ \hline
ВЫШЕ 2\NT=натурально, шлемовая рука, резкий расклад& ... \\ \hline
\end{tabularx}



\section{Ответы на открытие блоком}
\textbf{7. Ответы на 2\dia/2\hrt/2\spd}
\begin{itemize}[noitemsep]
    \item Фит на 3 уровне = 3+ карты, заграждение (НЕ ИНВИТ).
    \item Новая масть = 5+ карт, ФГ.
    \item 2\NT\ = ИНВИТ+ (запрос силы и концентрации).
\end{itemize}

\section{Важные штуки, перенести потом}
-октрытие--(вход)--\textbf{контра===инвит или ФГ без держки и фита, произвольный}

-открытие--(контра)--\textbf{реконтра===ОТ 10 очков}

-открытие--(вход)--\textbf{масть оппонента===ФГ с держкой или пониманием, что будет козырь}


-трансефры не сохраняются после перекрытия, НО

-(открытие)--1/2 БК--(пас)--\textbf{...сохранение торговли как после октрытие 1/2 бк}


-(1 минор)--контра--(пас)--\textbf{повтор минора===инвит+ без мажорных четверок}

-пас--(откртыие)--пас--(???)--\textbf{контра=== 10-11 очков, 4 карты во всех мастях}
\section{Конкурентная торговля (на открытия оппонентов)}
\subsection{Оппонент открылся на 1 ур., наш вход}
\begin{itemize}[label=\textcolor{sputnikRed}{\tiny\(\blacksquare\)}, noitemsep]
    \item 1\NT\ = 15-19 очков равномера с держкой.
    \item Масть без прыжка = (от 8 на первом уровне) 12-15 очков, 5+ карт.
    \item Масть с прыжком=слабая рука, блок.
    \item Контра =
    
     12-14 с равномером
    
    /15-17 очков равномера БЕЗ ДЕРЖКИ
    
    /16+очков с длинной мастью
    
    /18+ очков с равномером.
\end{itemize}

\subsection{Оппонент открылся блоком (2/3), наш вход}
\begin{itemize}[label=\textcolor{sputnikRed}{\tiny\(\blacksquare\)}, noitemsep]
    \item Масть = 12-17 очков, 5+ карт.
    \item Контра = 12+ очков (силовая/вызывающая).
    \item 2\NT/3\NT\ = 16-18 очков, задержка.
\end{itemize}

\section{Конкурентная торговля (на входы оппонентов)}
\subsection{Партнер открылся на 1, оппонент вошел}
\begin{itemize}[label=\textcolor{sputnikRed}{\tiny\(\blacksquare\)}, noitemsep]
    \item
\end{itemize}

\section{Вист}
\subsection{Первые ходы}
\begin{itemize}[noitemsep]
    \item Из двух карт: всегда старшая (\textbf{К}В, \textbf{Д}х, \textbf{8}2).
    \item Из ряда фигур: старшая (\textbf{Т}Кх, \textbf{К}Дх, \textbf{Д}Вх).
    \item Есть фигура (не подряд): 4-я сверху (\textbf{Т}76\textbf{3}, \textbf{К}В4\textbf{3}).
    \item Нет фигуры: 2-я сверху (\textbf{9}765, \textbf{10}73).
\end{itemize}
\subsection{Сигналы}
\begin{itemize}[noitemsep]
    \item \textbf{Маленькая} – поощрение (есть сила).
    \item \textbf{Старшая} (не фигура) – отказ.
    \item Снос фигуры обещает фигуру на 1 ранг ниже.
\end{itemize}

\end{document}
